% ==============================================================================
% Configurações Gerais do Template — Seminários em Engenharia Mecânica (DEM@UA)
% ==============================================================================

% --- Codificação e Idioma ---
\usepackage[utf8]{inputenc}
\usepackage[T1]{fontenc}
\usepackage[portuguese]{babel}
\usepackage{textcomp}
\usepackage{microtype}

% --- Tipografia: Apenas DUAS fontes ---
% 1) Corpo de texto (serif): Charter
% 2) Capa, títulos, cabeçalhos, rodapés e tabelas (sans-serif): Roboto
\usepackage{charter}
\usepackage{roboto}

% --- Margens e Espaçamento ---
\usepackage[
  a4paper,
  top=2.5cm,
  bottom=2.5cm,
  left=2.5cm,
  right=2.5cm,
  headheight=22pt
]{geometry}
\linespread{1.15}
\setlength{\parindent}{0pt}
\setlength{\parskip}{6pt}

% --- Paleta de Cores (Tons neutros + verde institucional UA/DEM, sem azulados) ---
\usepackage[table]{xcolor}
\definecolor{uagreen}{HTML}{6BA61E}      % Verde institucional DEM/UA
\definecolor{uadark}{HTML}{222222}       % Grafite neutro escuro para títulos e destaque
\definecolor{uamuted}{HTML}{595959}      % Cinzento neutro médio para rodapés/legendas
\definecolor{ualightbg}{HTML}{F7F7F7}    % Cinzento neutro muito claro para caixas
\definecolor{uatablehead}{HTML}{EBEBEB}  % Cinzento neutro suave para cabeçalhos de tabela
\definecolor{uaborder}{HTML}{D4D4D4}     % Cinzento neutro para linhas de separação

% --- Matemática e Números (separador decimal PT) ---
\usepackage{amsmath, amssymb}
\usepackage{icomma}

% --- Figuras, Tabelas e Listas ---
\usepackage{graphicx}
\graphicspath{{figs/}}
\usepackage{float}
\usepackage{booktabs}
\usepackage{tabularx}
\usepackage{array}
\usepackage{enumitem}
\usepackage[font={small,sf}, labelfont={bf,color=uadark}, skip=6pt]{caption}

% Espaçamento vertical confortável nas tabelas
\renewcommand{\arraystretch}{1.35}
\newcolumntype{L}[1]{>{\raggedright\arraybackslash}p{#1}}
\newcolumntype{R}[1]{>{\raggedleft\arraybackslash}p{#1}}
\newcolumntype{Y}{>{\raggedright\arraybackslash}X}

% --- Hierarquia de Títulos (Roboto, maiores do que o texto 11pt) ---
\usepackage{titlesec}

\titleformat{\section}
  {\sffamily\LARGE\bfseries\color{uadark}}
  {\textcolor{uagreen}{\thesection}}{0.6em}{}
  [\vspace{2pt}{\color{uagreen}\titlerule[1.2pt]}]
\titlespacing*{\section}{0pt}{22pt}{10pt}

\titleformat{\subsection}
  {\sffamily\Large\bfseries\color{uadark}}
  {\textcolor{uagreen}{\thesubsection}}{0.5em}{}
\titlespacing*{\subsection}{0pt}{16pt}{6pt}

\titleformat{\subsubsection}
  {\sffamily\large\bfseries\color{uamuted}}
  {\textcolor{uagreen}{\thesubsubsection}}{0.5em}{}
\titlespacing*{\subsubsection}{0pt}{12pt}{4pt}

% --- Cabeçalhos e Rodapés ---
\usepackage{fancyhdr}
\usepackage{lastpage}
\pagestyle{fancy}
\fancyhf{}
\fancyhead[L]{\sffamily\footnotesize\color{uamuted}Seminários em Engenharia Mecânica \textbar{} DEM -- UA}
\fancyhead[R]{\sffamily\footnotesize\color{uamuted}\nouppercase{\leftmark}}
\fancyfoot[L]{\sffamily\footnotesize\color{uamuted}Mestrado em Engenharia Mecânica 2026/2027}
\fancyfoot[R]{\sffamily\footnotesize\color{uadark}\textbf{\thepage} / \pageref{LastPage}}
\renewcommand{\headrulewidth}{0.4pt}
\renewcommand{\headrule}{\hbox to\headwidth{\color{uaborder}\leaders\hrule height \headrulewidth\hfill}}
\renewcommand{\footrulewidth}{0pt}

\fancypagestyle{frontmatter}{
  \fancyhf{}
  \fancyhead[L]{\sffamily\footnotesize\color{uamuted}Seminários em Engenharia Mecânica \textbar{} DEM -- UA}
  \fancyhead[R]{\sffamily\footnotesize\color{uamuted}\nouppercase{\leftmark}}
  \fancyfoot[L]{\sffamily\footnotesize\color{uamuted}Mestrado em Engenharia Mecânica 2026/2027}
  \fancyfoot[R]{\sffamily\footnotesize\color{uadark}\textbf{\thepage}}
  \renewcommand{\headrulewidth}{0.4pt}
  \renewcommand{\headrule}{\hbox to\headwidth{\color{uaborder}\leaders\hrule height \headrulewidth\hfill}}
  \renewcommand{\footrulewidth}{0pt}
}

% --- Caixa Moderna de Conteúdo (quebra automaticamente entre páginas) ---
\usepackage[most]{tcolorbox}
\newtcolorbox{campotexto}{
  breakable,
  enhanced,
  colback=ualightbg,
  colframe=ualightbg,
  borderline west={2.5pt}{0pt}{uagreen},
  boxrule=0pt,
  arc=2pt,
  left=10pt,
  right=10pt,
  top=8pt,
  bottom=8pt,
  before skip=6pt,
  after skip=10pt
}

% Estilo discreto para texto de orientação (substituir pelo conteúdo final)
\newcommand{\instrucao}[1]{\leavevmode{\color{uamuted}\itshape #1}}

% --- Bibliografia (BibTeX / natbib, estilo numérico) ---
\usepackage[numbers,sort&compress]{natbib}
\renewcommand{\bibsection}{}
\setlength{\bibsep}{4pt}

% --- Gráficos Vetoriais (para a Capa) ---
\usepackage{tikz}
\usetikzlibrary{calc}

% --- Hiperligações e Metadados PDF (sem cores azuladas) ---
\usepackage[
  colorlinks=true,
  linkcolor=uadark,
  urlcolor=uagreen,
  citecolor=uagreen,
  pdfborder={0 0 0}
]{hyperref}
\urlstyle{sf}

% --- Comandos para Metadados do Relatório ---
\newcommand{\tituloProjeto}{Título do Projeto de Inovação e I\&D}
\newcommand{\acronimoProjeto}{ACRÓNIMO}
\newcommand{\numeroGrupo}{Grupo \#}
\newcommand{\anoLetivo}{2026/2027}
